\clearpage\pagestyle{otMain}\setlength{\otBlocEspacement}{2.5mm}\fontsize{11}{15}\selectfont\corrigesChapitre{18}{Solides : cubes, pavés et patrons}
\begin{blocLivre}{cor-18-1}\begin{corrige}{1}{Nommer les solides}\demarche{Chercher les faces planes uniquement.}
Cube, pavé, cylindre, pyramide, cône, boule. Les polyèdres, aux faces toutes plates : $\res{\text{cube, pavé, pyramide}}$ - le cylindre, le cône et la boule ont des surfaces courbes.
\label{fig-cor-18-1}\figureGeo[corrige]{sol-1}
\end{corrige}\end{blocLivre}
\begin{blocLivre}{cor-18-2}\begin{corrige}{2}{Compter sur le cube}\demarche{Compter aussi les éléments cachés.}
Tous deux ont $\res6$ faces, $\res{12}$ arêtes et $\res8$ sommets : carrées pour le cube, rectangulaires pour le pavé.
\label{fig-cor-18-2}\figureGeo[corrige]{sol-2}
\end{corrige}\end{blocLivre}
\begin{blocLivre}{cor-18-3}\begin{corrige}{3}{Le prisme à base triangulaire}\demarche{Repérer d’abord les deux bases jumelles.}
Deux bases $\res{\text{triangulaires}}$ superposables, trois faces $\res{\text{rectangulaires}}$, 9 arêtes et 6 sommets.
\label{fig-cor-18-3}\figureGeo[corrige]{sol-3}
\end{corrige}\end{blocLivre}
\begin{blocLivre}{cor-18-4}\begin{corrige}{4}{La pyramide}\demarche{Séparer base et faces latérales.}
Une base carrée et 4 faces triangulaires : $\res5$ faces, $\res8$ arêtes et $\res5$ sommets, dont le sommet principal.
\label{fig-cor-18-4}\figureGeo[corrige]{sol-4}
\end{corrige}\end{blocLivre}
\begin{blocLivre}{cor-18-5}\begin{corrige}{5}{Patron de cube ?}\demarche{Plier mentalement face après face.}
La croix et la forme en T se replient sans superposition : $\res{\text{patrons}}$. La disposition contenant un bloc carré de quatre cases superpose deux faces : $\res{\text{pas un patron}}$.
\label{fig-cor-18-5}\figureGeo[corrige]{sol-5}
\end{corrige}\end{blocLivre}
\begin{blocLivre}{cor-18-6}\begin{corrige}{6}{Les faces opposées}\demarche{Sauter une face dans l’alignement.}
Les faces opposées sont séparées par une face sur le patron : l’opposée de A est $\res C$, celle de E est $\res F$.
\label{fig-cor-18-6}\figureGeo[corrige]{sol-6}
\end{corrige}\end{blocLivre}
\begin{blocLivre}{cor-18-7}\begin{corrige}{7}{Patron de pavé ?}\demarche{Comparer les bords qui se rejoignent.}
La disposition 1 réunit bien $\res{\text{trois paires}}$ de rectangles identiques, bords communs de même longueur : c’est un patron. Dans la 2, deux bords pliés ensemble ont des longueurs différentes : le pliage bâille.
\label{fig-cor-18-7}\figureGeo[corrige]{sol-7}
\end{corrige}\end{blocLivre}
\begin{blocLivre}{cor-18-8}\begin{corrige}{8}{Le patron du dé}\demarche{Former les paires qui totalisent 7.}
L’opposée de 3 porte $\res4$, celle de 2 porte $\res5$, celle de 1 porte $\res6$ : les valeurs, en vert, respectent les paires séparées d’une face.
\label{fig-cor-18-8}\figureGeo[corrige]{sol-8}
\end{corrige}\end{blocLivre}
\begin{blocLivre}{cor-18-9}\begin{corrige}{9}{Construire le patron}\demarche{Prévoir les trois paires avant de tracer.}
Trois paires : deux faces de $4\times2$, deux de $4\times1$, deux de $2\times1$. Le patron, en vert à échelle réduite, aligne ces six rectangles en respectant les bords communs.
\label{fig-cor-18-9}\figureGeo[corrige]{sol-9}
\end{corrige}\end{blocLivre}
\begin{blocLivre}{cor-18-10}\begin{corrige}{10}{Les cubes cachés}\demarche{Compter étage par étage.}
Étage du bas : $4\times2=8$ ; étage du milieu : 4 ; étage du haut : 2. Total : $8+4+2=\res{14}$ cubes, porteurs cachés comptés.
\label{fig-cor-18-10}\figureGeo[corrige]{sol-10}
\end{corrige}\end{blocLivre}
\begin{blocLivre}{cor-18-11}\begin{corrige}{11}{La bonne vue}\demarche{Imaginer chaque point de vue séparément.}
La vue de face lit les hauteurs des colonnes : vignette $\res2$. La vue de dessus lit les cases occupées au sol : vignette $\res1$. La vignette 3 ne correspond à aucun point de vue.
\label{fig-cor-18-11}\figureGeo[corrige]{sol-11}
\end{corrige}\end{blocLivre}
\begin{blocLivre}{cor-18-12}\begin{corrige}{12}{Compléter le grand cube}\demarche{Calculer le volume visé en cubes.}
Le grand cube contient $3\times3\times3=27$ cubes : il manque $27-20=\res7$ cubes.
\label{fig-cor-18-12}\figureGeo[corrige]{sol-12}
\end{corrige}\end{blocLivre}
\begin{blocLivre}{cor-18-13}\begin{corrige}{13}{Coder le chemin}\demarche{Mettre à jour l’orientation à chaque virage.}
Un codage correct : avancer de 3 ; pivoter à gauche ; avancer de 2 ; pivoter à droite ; avancer de 2 ; pivoter à droite ; avancer de 2. Les quarts de tour suivent l’orientation $\res{\text{du pion}}$, pas celle du lecteur.
\label{fig-cor-18-13}\figureGeo[corrige]{sol-13}
\end{corrige}\end{blocLivre}
\begin{blocLivre}{cor-18-14}\begin{corrige}{14}{Exécuter le programme}\demarche{Suivre le rectangle case par case.}
Le chemin dessine un rectangle de 4 cases sur 2 : le pion revient $\res{\text{à sa case de départ}}$, désormais tourné vers le bas de la feuille. Contrôle : deux côtés de 4, deux côtés de 2, trois quarts de tour à gauche.
\label{fig-cor-18-14}\figureGeo[corrige]{sol-14}
\end{corrige}\end{blocLivre}
\begin{blocLivre}{cor-18-15}\begin{corrige}{15}{Le cube peint}\demarche{Classer les cubes selon leur position.}
Deux faces peintes : les cubes du milieu des arêtes, $\res{12}$. Aucune : le cube central, $\res1$. Contrôle : $8$ coins $+12+6$ centres de faces $+1=27$.
\label{fig-cor-18-15}\figureGeo[corrige]{sol-15}
\end{corrige}\end{blocLivre}
